\begin{figure}[htbp]
\centering
\begin{tikzpicture}[x=1cm,y=1cm]
\node[booknode,text width=5.4cm] (p) at (6,0) {父操作：账户、购买动作、K1\\服务保存稳定购买编号 P1};
\node[booknode,text width=3.5cm] (d) at (1.5,-2) {扣款子动作\\(P1, debit, 1)\\扣款重试继续此身份};
\node[booknode,text width=3.5cm] (n) at (6,-2) {通知子动作\\(P1, notice, 1)\\通知重试继续此身份};
\node[booknode,text width=3.5cm] (i) at (10.5,-2) {发票子动作\\(P1, invoice, 1)\\发票重试继续此身份};
\draw[bookarrow] (p)--(d);
\draw[bookarrow] (p)--(n);
\draw[bookarrow] (p)--(i);
\node[booknote,text width=12cm] at (6,-3.6)
 {接收服务再加自己的已认证调用方与动作范围；箭头说明来源和调用关系，跨服务共同提交另需协议};
\end{tikzpicture}
\caption{父操作相同，子动作仍需分别识别；同一子动作重试继续沿用自己的身份。}
\label{fig:ch03-derived-identity}
\end{figure}
